% GENERATED FILE. Edit canonical lessons, then run npm run generate:books.
% canonical-source-hash: fnv1a64:65d4d81eaab31842
% canonical-lessons: BN-C72-boro, BN-C72-chhoto, BN-C72-notun, BN-C72-purono, BN-C72-kharap

\chapter{Big, Small, New, Old, Bad}
\label{ch:bn-big-small-new-old-bad}
\hlchaptermodality{72}

\begin{chapteropening}
\textbf{By the end of this chapter:} \emph{I can say big, small, new, old and bad.}

Complete the last lesson of chapter 72: I can say big, small, new, old and bad.
\end{chapteropening}

\section[\texorpdfstring{\bn{বড়} (bôṛo)}{বড় (bôṛo)}]{\bn{বড়} (bôṛo) — big}

\label{lesson:BN-C72-boro}



\begin{quote}
Before the new one: say the Bengali for an ant, then the Bengali for a hen.
\end{quote}

\subsection*{You'll want to know: \bn{বড়}}
\textbf{\bn{বড়}} — \emph{bôṛo} — \textquotedblleft{}big\textquotedblright{}.

\begin{cousinweb}
\textbf{\bn{বড়}} (\emph{bôṛo}) is a tadbhava, a word inherited from Sanskrit and reshaped by sound change. It comes from \emph{vaḍra}, \textquotedblleft{}big\textquotedblright{} (CDIAL 11225). CDIAL rejects, as unlikely, a derivation from \emph{vṛddha}, \textquotedblleft{}grown\textquotedblright{}.
\end{cousinweb}

\begin{grammarlens}[title={What you've built}]
One more word for this chapter.
\end{grammarlens}

\subsection*{Your turn}
\begin{itemize}
\raggedright
  \item \emph{Say it:} \emph{bôṛo}
  \item \emph{Say it:} \emph{bôṛo}, once more
  \item \emph{Say it:} say \emph{bôṛo}
  \item \emph{Recall:} say the Bengali for an ant, then the Bengali for a hen, then say \emph{bôṛo} again
\end{itemize}

\begin{tcolorbox}[breakable,colback=teal!4,colframe=teal!35!black,title={Before you move on}]
What does \bn{বড়} mean? (\textquotedblleft{}Big\textquotedblright{}.) Say it once more.
\end{tcolorbox}
\section[\texorpdfstring{\bn{ছোট} (chhoṭo)}{ছোট (chhoṭo)}]{\bn{ছোট} (chhoṭo) — small}

\label{lesson:BN-C72-chhoto}



\begin{quote}
Before the new one: say the Bengali for a hen, then the Bengali for big.
\end{quote}

\subsection*{You'll want to know: \bn{ছোট}}
\textbf{\bn{ছোট}} — \emph{chhoṭo} — \textquotedblleft{}small\textquotedblright{}.

\begin{cousinweb}
\textbf{\bn{ছোট}} (\emph{chhoṭo}) is a tadbhava, a word inherited and reshaped by sound change. It comes from \emph{*choṭṭa}, \textquotedblleft{}small\textquotedblright{}, a form CDIAL reconstructs from its descendants (CDIAL 5071).
\end{cousinweb}

\begin{grammarlens}[title={What you've built}]
One more word for this chapter.
\end{grammarlens}

\subsection*{Your turn}
\begin{itemize}
\raggedright
  \item \emph{Say it:} \emph{chhoṭo}
  \item \emph{Say it:} \emph{chhoṭo}, once more
  \item \emph{Say it:} say \emph{chhoṭo}
  \item \emph{Recall:} say the Bengali for a hen, then the Bengali for big, then say \emph{chhoṭo} again
\end{itemize}

\begin{tcolorbox}[breakable,colback=teal!4,colframe=teal!35!black,title={Before you move on}]
What does \bn{ছোট} mean? (\textquotedblleft{}Small\textquotedblright{}.) Say it once more.
\end{tcolorbox}
\section[\texorpdfstring{\bn{নতুন} (nôtun)}{নতুন (nôtun)}]{\bn{নতুন} (nôtun) — new}

\label{lesson:BN-C72-notun}



\begin{quote}
Before the new one: say the Bengali for big, then the Bengali for small.
\end{quote}

\subsection*{You'll want to know: \bn{নতুন}}
\textbf{\bn{নতুন}} — \emph{nôtun} — \textquotedblleft{}new\textquotedblright{}.

\begin{cousinweb}
\textbf{\bn{নতুন}} (\emph{nôtun}) is usually traced to Sanskrit \emph{nūtana}, \textquotedblleft{}new, fresh\textquotedblright{}, with its vowels rearranged. CDIAL has that entry but lists no Bengali form under it (CDIAL 7575).
\end{cousinweb}

\begin{grammarlens}[title={What you've built}]
One more word for this chapter.
\end{grammarlens}

\subsection*{Your turn}
\begin{itemize}
\raggedright
  \item \emph{Say it:} \emph{nôtun}
  \item \emph{Say it:} \emph{nôtun}, once more
  \item \emph{Say it:} say \emph{nôtun}
  \item \emph{Recall:} say the Bengali for big, then the Bengali for small, then say \emph{nôtun} again
\end{itemize}

\begin{tcolorbox}[breakable,colback=teal!4,colframe=teal!35!black,title={Before you move on}]
What does \bn{নতুন} mean? (\textquotedblleft{}New\textquotedblright{}.) Say it once more.
\end{tcolorbox}
\section[\texorpdfstring{\bn{পুরনো} (purôno)}{পুরনো (purôno)}]{\bn{পুরনো} (purôno) — old}

\label{lesson:BN-C72-purono}



\begin{quote}
Before the new one: say the Bengali for small, then the Bengali for new.
\end{quote}

\subsection*{You'll want to know: \bn{পুরনো}}
\textbf{\bn{পুরনো}} — \emph{purôno} — \textquotedblleft{}old\textquotedblright{}.

\begin{cousinweb}
\textbf{\bn{পুরনো}} (\emph{purôno}) is a tadbhava, a word inherited from Sanskrit and reshaped by sound change. It comes from \emph{purāṇa}, \textquotedblleft{}ancient\textquotedblright{} (CDIAL 8283).
\end{cousinweb}

\begin{grammarlens}[title={What you've built}]
One more word for this chapter.
\end{grammarlens}

\subsection*{Your turn}
\begin{itemize}
\raggedright
  \item \emph{Say it:} \emph{purôno}
  \item \emph{Say it:} \emph{purôno}, once more
  \item \emph{Say it:} say \emph{purôno}
  \item \emph{Recall:} say the Bengali for small, then the Bengali for new, then say \emph{purôno} again
\end{itemize}

\begin{tcolorbox}[breakable,colback=teal!4,colframe=teal!35!black,title={Before you move on}]
What does \bn{পুরনো} mean? (\textquotedblleft{}Old\textquotedblright{}.) Say it once more.
\end{tcolorbox}
\section[\texorpdfstring{\bn{খারাপ} (khārāp)}{খারাপ (khārāp)}]{\bn{খারাপ} (khārāp) — bad}

\label{lesson:BN-C72-kharap}



\begin{quote}
Before the new one: say the Bengali for new, then the Bengali for old.
\end{quote}

\subsection*{You'll want to know: \bn{খারাপ}}
\textbf{\bn{খারাপ}} — \emph{khārāp} — \textquotedblleft{}bad\textquotedblright{}.

\begin{cousinweb}
\textbf{\bn{খারাপ}} (\emph{khārāp}) is a loan from Arabic \emph{kharāb}, \textquotedblleft{}ruin, desolation\textquotedblright{}, which reached Bengali through Persian. Its Arabic root is \emph{kh-r-b} (Lane). In Bengali it means \textquotedblleft{}bad\textquotedblright{}.
\end{cousinweb}

\begin{grammarlens}[title={What you've built}]
One more word for this chapter.
\end{grammarlens}

\subsection*{Your turn}
\begin{itemize}
\raggedright
  \item \emph{Say it:} \emph{khārāp}
  \item \emph{Say it:} \emph{khārāp}, once more
  \item \emph{Say it:} say \emph{khārāp}
  \item \emph{Recall:} say the Bengali for new, then the Bengali for old, then say \emph{khārāp} again
\end{itemize}

\begin{tcolorbox}[breakable,colback=teal!4,colframe=teal!35!black,title={Before you move on}]
What does \bn{খারাপ} mean? (\textquotedblleft{}Bad\textquotedblright{}.) Say it once more.
\end{tcolorbox}
